\documentclass{article}

\usepackage{amsmath,amssymb}
\usepackage{xurl}
\usepackage[hidelinks]{hyperref}
\setlength{\emergencystretch}{2em}

% Shared commands for notes in docs/paper-gaps/.
% Notation follows the LDT paper (references/ldt-paper/).

\newcommand{\Lean}{\textsc{Lean}}
\newcommand{\leanid}[1]{\textsf{#1}}
\newcommand{\ghissue}[1]{\href{https://github.com/LionSR/MIPStarRE/issues/#1}{GitHub issue \##1}}
\newcommand{\ghpr}[1]{\href{https://github.com/LionSR/MIPStarRE/pull/#1}{GitHub PR \##1}}

% --- Paper-compatible notation ---
\newcommand{\F}{\mathbb{F}}
\newcommand{\N}{\mathbb{N}}
\newcommand{\C}{\mathbb{C}}
\newcommand{\tr}{\operatorname{tr}}
\newcommand{\ot}{\otimes}
\newcommand{\E}{\mathop{\bf E\/}}
\newcommand{\bx}{\mathbf{x}}
\newcommand{\by}{\mathbf{y}}
\newcommand{\bu}{\mathbf{u}}
\newcommand{\bv}{\mathbf{v}}

% Approximation relation (paper uses \approx_\eps and \simeq_\zeta)
\newcommand{\approxeps}{\approx_{\varepsilon}}
\newcommand{\simgeq}{\simeq_{\zeta}}

% Polynomial spaces (paper uses \polyfunc, \polysub, \polymeas)
\newcommand{\polyfunc}[3]{\operatorname{Poly}(#1,#2,#3)}
\newcommand{\polysub}[3]{\operatorname{PolySub}(#1,#2,#3)}
\newcommand{\polymeas}[3]{\operatorname{PolyMeas}(#1,#2,#3)}


\title{Resolved Missing $\nu$-Consistency Hypothesis in the Self-Improvement Theorem}
\author{}
\date{2026-05-01}

\begin{document}
\maketitle

\section{At a glance}

\begin{itemize}
  \item \textbf{Status: resolved.}  The statement-level discrepancy is resolved:
  the Lean theorem \leanid{selfImprovement} contains the paper's
  $\nu$-consistency hypothesis.  The proof-level gap described in the original
  note has also been closed: \leanid{selfImprovement} and the induction-section
  reformulation are now proved from the corrected source-facing hypotheses.
  Thus issue~\#1515 is historical for the current Lean tree.
  \item \textbf{Difficulty: low as a statement repair, substantial as the
  former proof construction.}  The missing hypothesis was easy to identify from
  the paper statement.  The later proof closure required the Section~9 helper,
  semidefinite-programming, orthonormalization, and final-field constructions.
  \item \textbf{Estimated weight:} no remaining proof construction is expected
  for this note.  It remains as documentation of the former statement-level
  discrepancy and its resolution.
  \item \textbf{Mathlib/project split:} roughly \textbf{25\% Mathlib} and
  \textbf{75\% project-local} for the proof route that closed the issue.
  Mathlib supplies finite-dimensional operator and matrix order infrastructure;
  the project supplies the LDT self-improvement interfaces and transports.
  \item \textbf{Key Mathlib inputs:} positive semidefinite matrices, finite
  sums of matrix-valued effects, trace and operator inequalities, and the
  finite-dimensional SDP infrastructure used by the Section~9 proof.
\end{itemize}

\section{Key theorem forms}

The self-improvement step in \cite[Section~9]{Ji2020LowIndividualDegree}
takes a measurement $G\in\polymeas{m}{q}{d}$, where $m$ is the number
of variables, $q$ the field size, and $d$ the degree bound of the
low individual degree test.  The surrounding strategy is an
$(\eps,\delta,\gamma)$-good symmetric strategy; its data includes a
projective measurement family $A=\{A^{\bu}\}_{\bu\in\F_q^{\,m}}$ with
outcomes in $\F_q$ and a bipartite state $\ket\psi$.  The function $g$
maps a point $\bu$ to its evaluation, and $[g(\bu)=a]$ selects the outcome
indexed by $a$.  The paper's self-improvement theorem then assumes the
following consistency hypothesis on $G$: on average over $\bu\in\F_q^{\,m}$,
\begin{equation}
  \label{eq:paper-nu-consistency}
  A_a^{\bu}\ot I \simeq_{\nu} I \ot G_{[g(\bu)=a]}.
\end{equation}
Here $\simeq_{\nu}$ is the paper's approximation relation comparing two
families of measurements by the expectation of their difference: for a
bipartite state $\ket\psi$, the expectation difference is at most $\nu$.
The theorem produces an output measurement $H$ and a dual certificate
operator $Z$ dominating the averaged $A$-operators.
The parameter $\zeta$ is set to
$\zeta = 3000\,m\,(\eps^{1/32}+\delta^{1/32}+(d/q)^{1/32})$.
Under \eqref{eq:paper-nu-consistency} the theorem concludes completeness
$\bra\psi H\ot I\ket\psi \ge (1-\nu)-\zeta$, point consistency of $H$,
strong self-consistency of $H$ at level $\zeta$, and the operator bound
$\bra\psi Z\ot(I-H)\ket\psi \le \zeta$.\footnote{The paper passages are
\path{references/ldt-paper/self_improvement.tex}, lines~635--671 (theorem)
and lines~24--60 (helper lemma); the corresponding induction-section theorem is in
\path{references/ldt-paper/inductive_step.tex}, lines~249--286. This note
is part of the retrospective audit in \ghissue{930}. The discrepancy is
tracked for resolution in \ghissue{931}.}

\section{The formal statement}

The Lean declaration \leanid{MIPStarRE.LDT.SelfImprovement.selfImprovement}
now includes the consistency hypothesis
\eqref{eq:paper-nu-consistency}.  The former obligation bundle
\leanid{SelfImprovementObligations} and the conditional theorem that used it
have been removed.  Thus the discrepancy recorded in this note has been
resolved at the level of the public theorem statement.

The proof-level gap recorded by the original version of this note has since
been discharged.  The current theorem proves the source-shaped statement
directly, rather than representing the missing work by an explicit
\leanid{sorry} or by extra hypotheses on a conditional substitute for the
theorem.

The induction-step theorem
\leanid{MIPStarRE.LDT.MainInductionStep.selfImprovementInInductionSection}
has the same status: its statement includes the input consistency hypothesis,
and the Section~9 work needed by that reformulation is now checked.

\section{Mathematical consequence}

The previous theorem statement was too weak to express the paper's claim,
because $\nu$ was not related to the input measurement.  The present theorem
statement no longer has that defect.  The helper strong self-consistency
estimate, the orthonormalization step, and the final-field transport are now
derived inside the checked proof rather than supplied as auxiliary source
hypotheses.

\section{Blueprint status}

The blueprint entry now points to the theorem \leanid{selfImprovement}.  No
top-level conditional obligation theorem is recorded as the formalization of
the paper theorem.  The former proof-level gap tracked by issue~\#1515 is
historical for the current code.

\section{Error coefficients}

The paper's helper lemma carries its own error parameter
$\widehat\zeta = 100\,m\,(\eps^{1/2}+\delta^{1/2}+(d/q)^{1/2})$;
the orthonormalization step contributes a further error
$\widehat\zeta_{\text{ortho}} = 100\,\widehat\zeta^{\,1/4}$.
All error constants match the paper when the obligations are discharged.
The helper error matches
$\widehat\zeta$; the orthogonalization error delegates to the
orthonormalization infrastructure, which uses the paper's $100\widehat\zeta^{\,1/4}$;
the data-processing error $8\widehat\zeta+8\sqrt{\widehat\zeta_{\text{ortho}}}$
matches the paper; and the final self-improvement error
$\zeta = 3000\,m\,(\eps^{1/32}+\delta^{1/32}+(d/q)^{1/32})$ matches the paper.

\section{Conclusion}

The paper's self-improvement theorem requires an approximation hypothesis
$\simeq_{\nu}$ linking the input measurement $G$ to the parent measurement
$A$.  The Lean statement now contains this hypothesis, and the corresponding
proof is checked without reintroducing the missing work as auxiliary
hypotheses.

\bibliographystyle{alpha}
\bibliography{references}

\end{document}
